\hyperref[def:exchange-rate-regimes-capital-flows]{Exchange-rate regimes and their effects} follows the quote convention in \hyperref[def:fx-market-real-exchange-rates]{the FX market}.

\begin{qfdefinition}[Exchange-rate regimes and their effects: the ideal regime, the spectrum from dollarisation to free floating, effects on trade through elasticities and absorption, capital flows, and the objectives of capital restrictions]\label{def:exchange-rate-regimes-capital-flows}
\qffield{Statement} \emph{The ideal regime} would combine a credibly fixed rate, full convertibility and an independent monetary policy, which the impossible trinity rules out \citep[chap.~9, sec.~4.1, p.~500]{economics-investment-decision-makers-2013}. \emph{The spectrum} (sec.~4.3, pp.~503--511), from least to most flexible:
\begin{itemize}
\item \emph{arrangements with no separate legal tender}: dollarisation, or a monetary union, which gives up independent policy;
\item a \emph{currency board}: domestic currency fully backed by reserves of the anchor at a fixed rate, with monetary policy surrendered;
\item a \emph{fixed parity} (conventional peg) within narrow margins, defended by intervention;
\item \emph{target zones} (wider bands) and \emph{crawling pegs} adjusted periodically, often for inflation differentials, and crawling bands;
\item a \emph{managed float}, intervening without a stated target;
\item an \emph{independent float}, the rate set by the market.
\end{itemize}
More fixed regimes import credibility but leave the economy exposed to speculative attack if reserves or commitment are doubted; more flexible regimes absorb shocks but add volatility.

\emph{Effects on trade} (sec.~5, pp.~512--520). In the \emph{elasticities approach} a depreciation improves the trade balance if the generalised Marshall--Lerner condition holds,
\[
\omega_X\varepsilon_X+\omega_M(\varepsilon_M-1)>0,
\]
reducing to $\varepsilon_X+\varepsilon_M>1$ when trade is balanced; demand responds with lags, so the balance may first worsen (the J-curve). The \emph{absorption approach} adds that the balance improves only if national income rises relative to spending, so at full employment a depreciation must be accompanied by lower absorption. \emph{Capital flows}: exchange-rate expectations and interest differentials drive flows, which in floating regimes often dominate short-run moves.

\emph{Capital restrictions} (chap.~8, sec.~3.5, pp.~435--437) aim to reduce volatility from hot money, defend a peg or allow independent monetary policy, keep domestic saving at home, and protect strategic industries from foreign ownership; they may also reduce efficiency and invite evasion.

\qffield{Notation} $\omega_X,\omega_M$ shares of exports and imports in total trade; $\varepsilon_X,\varepsilon_M$ price elasticities of foreign demand for exports and domestic demand for imports.

\qffield{Conditions} EIDM's taxonomy follows the IMF's de jure classification as described in 2013; de facto classifications differ (\hyperref[def:currency-regimes-intervention]{currency regimes}). Marshall--Lerner assumes exports priced in domestic and imports in foreign currency.

\qffield{Practice} \emph{Reproduced} (Exhibit~9-9, p.~514): a 10 percent euro depreciation raises import spending on EUR 600 million by 3.5 percent ($\varepsilon_M=0.65$), EUR 21 million, and export receipts on EUR 400 million by 7.5 percent ($\varepsilon_X=0.75$), EUR 30 million: the balance improves by EUR 9 million, and $0.4\times0.75+0.6\times(0.65-1)=0.09>0$. The book's printed working reads ``0.65 -- 4'' for $0.65-1$; the result 0.09 is correct.

\qffield{Pitfalls} Expecting a depreciation to improve the balance immediately. Treating a peg as a free option for investors: when it breaks, the move is large. Assuming capital controls work equally against persistent and temporary flows.

\qffield{See also} \hyperref[def:currency-regimes-intervention]{currency regimes and intervention}; \hyperref[def:flows-policy-exchange-rates]{payment flows and policy}; \hyperref[meth:trade-benefits-restrictions-blocs]{trade restrictions}.

\end{qfdefinition}

\begin{qfdefinition}[Currency regimes and intervention: fixed, managed and floating regimes in Gliner's grouping, sterilised and unsterilised intervention, the impossible trinity, and reserve depletion]\label{def:currency-regimes-intervention}
\qffield{Statement} Gliner groups currencies by the central bank's intervention: \emph{fixed} (HKD, DKK, LVL, AED, TRL in his 2013 table), \emph{floating with frequent intervention} (JPY, BRL, CNY, CHF with the EUR/CHF floor) and \emph{floating with no intervention} (USD, EUR, GBP, AUD, NZD, MXN, CAD, SEK, NOK, INR). He stresses that ``no model will be of help if, for example, a currency is pegged'' \citep[chap.~7, sec.\ ``Currency Regimes'', PDF pp.~185--194]{s7-gliner}. His examples include the Hong Kong dollar's peg at 7.8 (p.~190) and Brazil's repeated interventions (Table~7.2, pp.~191--192).

Intervention is \emph{sterilised} when the central bank offsets its effect on the monetary base and \emph{unsterilised} when it lets the base change (chap.~11, pp.~360--361). The \emph{impossible trinity}, from Mundell's work, says a country can have only two of a fixed exchange rate, free capital flows and an independent monetary policy (chap.~11, sec.\ ``The Impossible Trinity'', p.~363).

\qffield{Notation} Peg, crawling peg, band and currency board are forms of a fixed regime; a floor (EUR/CHF 1.20 from 6 September 2011) limits movement in one direction.

\qffield{Conditions} Gliner's grouping is de jure and as of 2013; de facto classifications can differ, and several of his currencies have since changed regime (the EUR/CHF floor ended in 2015).

\emph{A de facto check} uses the Ilzetzki--Reinhart--Rogoff monthly fine classification for December 2013 and December 2025 \citep[sheets ``Classification'' and ``Fine'', last update 5 October 2026]{irr-era-2025}. Its codes are 1 no separate legal tender, 2 pre-announced peg or currency board, 3 pre-announced horizontal band within $\pm2\%$, 11 moving band within $\pm2\%$, 12 de facto moving band within $\pm5\%$ or managed floating, and 13 freely floating. HKD, DKK and AED are 2 in both months, and Latvia moves from 2 to 1 in January 2014, when it adopted the euro. But Gliner's ``fixed'' TRL is 12 in both, and his intervening floats differ: CNY is 3 (a band) and CHF 2 in 2013 (the floor), 11 in 2025; JPY moves from 13 to 12 in August 2024.

Of his ``no intervention'' floats only USD and AUD are 13 in both months; GBP, CAD and NZD are 11 in 2025 (since January 2020), SEK and INR 11 in both, MXN and NOK reach 13 only in 2016 and 2023. So de jure ``free float'' and de facto ``freely floating'' are different lists.

\qffield{Practice} Standard arguments, not shown in Gliner's chapter, interpret the regimes: in a free float price discovery happens in FX; in a managed regime pressure appears in rates, reserves or capital controls; a hard peg looks stable until it breaks; and ``cheap'' can stay cheap if the regime prevents adjustment. The same applies to the usual intervention motives and conditions for success (reserve adequacy, credibility, consistency with monetary policy, liquidity), to the view that intervention against fundamentals typically fails, and to reserve depletion as a leading indicator. The trinity explains the ERM crisis of 1992 and the Asian crisis of 1997 in telling; Gliner gives no such cases in that section.

\qffield{Pitfalls} The trinity is about full capital mobility and a fixed rate; intermediate regimes trade off partial independence against partial stability, so ``choose two of three'' is a corner-solution statement.

\qffield{See also} \hyperref[def:central-bank-policy-framework]{central-bank tools}; \hyperref[rule:major-central-banks-current]{major central banks}; \hyperref[def:us-dollar-dxy-sdr]{the dollar and the SDR}.

\qffield{Further reading} On whether intervention works, \citep{cur-fratzscher-et-al-2019}; for a survey, \citep{cur-menkhoff-2013}.

\end{qfdefinition}
